\begin{lemma}[Existenz einer tubularen Umgebung\label{jordan:lemma:tubulare_umgebung}]
    Sei $\gamma$ eine reell-analytische Kurve in der Ebene $\mathbb{R}^2$. Dann existiert $\varepsilon > 0$, so dass die Abbildung $\varphi: S\textsuperscript{1} \times (-\varepsilon, \varepsilon) \to \mathbb{R}^2$, $\varphi(t,s) = \gamma(t) + s \cdot n(t)$, ein reell-analytischer Diffeomorphismus auf eine offene Umgebung $T$ von $\gamma$ ist, wobei $n(t)$ der nach aussen gerichtete Einheitsnormalenvektor ist. 
    Das Komplement $T \setminus \gamma$ zerfällt in zwei zusammenhängende Komponenten entsprechend $s > 0$ und $s < 0$, die als innere und äussere tubulare Umgebung bezeichnet werden.    
\end{lemma}